1. Purpose and Evaluation Context
This output is a graded academic artifact, not a creative task. The lecture scribe must be suitable for reading-based / closed-notes exam preparation and serve as faithful reference material for revising what was taught in class.
This task evaluates: Precision in prompt adherence; Faithfulness to provided material; Disciplined and reproducible use of AI.
This task does not evaluate: Creativity; Personal writing style; Fluency or verbosity beyond necessity.
The scribe must answer the question: “If a student studies only this document, can they reliably revise what was taught?”

2. Definition of a Lecture Scribe
A lecture scribe is a faithful, exam-ready reconstruction of lecture content. A correct lecture scribe must allow a student to: Recall exact definitions and notation; Reconstruct proofs and derivations; Follow worked examples step by step; Understand logical dependencies between concepts.

3. Context Constraint (Authoritative Sources Only)
The only allowed context consists of: The attached Lecture slides / lecture PDF; Attached textbook excerpts that are explicitly relevant to this lecture. You must not use any other source of information.
Not Allowed: Solved material not shown or discussed in lecture; Instructor meta-instructions on how to write a scribe; Reference scribes or ideal answers; Writing guides or explanations disguised as content; Intuition, commentary, or simplifications not present in lecture.
Everything included in the scribe must be directly traceable to the provided context.

4. Content Requirements
The lecture scribe must include, exactly as presented in the lecture: Definitions and notation; Assumptions and conditions; Statements of theorems or results; Proofs or proof sketches, reconstructed step by step; Worked examples only if they appear in the context, with all intermediate steps; Clear logical flow following the lecture’s structure and order.
You must not: Invent definitions, results, notation, or examples; Add explanations or intuition not explicitly discussed; Reorder topics beyond the lecture’s flow.

5. Reasoning Requirement (Chain-of-Thought)
You must include explicit reasoning steps in the scribe wherever the lecture includes: Proofs; Derivations; Example solutions.
Each reasoning step must follow only from: Previously stated definitions; Assumptions or conditions; Earlier results from the same lecture.
Reasoning must reflect how conclusions follow, not external intuition.

6. Structure Requirement
Organize the scribe using clear sectioning that mirrors the lecture structure (for example): Random Variables; Bernoulli Random Variable; Binomial Distribution; Geometric Distribution; Poisson Distribution.
Logical dependencies between concepts must be explicit.

7. Output Constraint
Output only the lecture scribe. No commentary before or after the scribe. No meta explanations about the process.

8. Self-Consistency Prompting Requirement
This prompt is designed for self-consistency evaluation. You must assume that: The same prompt will be reused across multiple independent runs; The same context will be reused across all runs.
Your task in each run is to generate one complete lecture scribe, fully self-contained. No reference to other runs is allowed.

9. Instruction to Begin
Generate the lecture scribe now, strictly following all constraints above.